\section{Proof of Theorem~\ref{thm:power saving}}

\subsection{Construction of eligible hypergraphs}\label{subsec:construction of eligible hypergraphs}

As discussed in the proof overview, our aim is to construct a sequence of hypergraphs $F_0, \dots, F_\ell$ satisfying certain properties. In this section we will state these properties precisely. In particular, we will show how to construct $F_{j+1}$ from $F_j$ while making sure that all the necessary properties are maintained. To be able to state these properties, we need to introduce a special kind of independent set in a hypergraph, which is particularly suited for gluing along it.

\begin{definition}\label{defn:good}
Let $F$ be a 3-uniform hypergraph. We say that an independent set $A\subseteq V(F)$ is \textit{good} if for any set $U\supsetneq A$, we have $\Delta(U)\geq |A|+1$.
\end{definition}

\noindent
Note that the collection of good sets is closed under taking subsets. To see this, suppose that we have a good set $A\subset V(F)$ and $A'\subset A$. To show that $A'$ is good, we need to verify that for any $U\supsetneq A'$ one has $\Delta(U)\geq |A'|+1$. If $U\subset A$, this is clear since $U$ must be an independent set and hence $\Delta(U)=|U|\geq |A'|+1$. Otherwise, by using that $A$ is good and that $A\subsetneq U\cup A$ we obtain $\Delta(U\cup A)\geq |A|+1$. This allows us to show the desired bound on $\Delta(U)$ as follows.
\begin{align*}
    \Delta(U)=v(F[U])-e(F[U])&\geq v(F[U\cup A])-|A\backslash U|-e(F[U\cup A])\\
    &\geq \Delta(U\cup A)-|A\backslash A'|\geq |A|-|A\backslash A'|+1=|A'|+1.
\end{align*}

This discussion allows us to conclude that for any good set $A$ and any set $U\subseteq V(F)$ not fully contained in $A$, the set $A\cap U$ is also a good set and therefore $\Delta(U)\geq |A\cap U|+1$. 

Our definition of good sets vaguely resembles Definition 2.3 in the paper of Conlon, Gishboliner, Levanzov and Shapira \cite{CGLS23}. Namely, they consider independent sets $A$ of size $\Delta(F)+1$ with the property that the deficiency of any set $U\subset V(F)$ is roughly controlled by $|A\cap U|$. However, their definition is more complex than ours, because their gluing procedure (relying on the hypergraph regularity lemma) glues several copies of the hypergraph at a time instead of $2$ as in our paper.

\begin{definition}\label{defn:eligible}
A 3-uniform hypergraph $F$ with $\Delta(F)=k\geq 1$ is \textit{eligible} if it satisfies the following properties:
\begin{itemize}
    \item[(i)] there are two disjoint good sets $A, B\subset V(F)$ of size $|A|=|B|=k-1$,
    \item[(ii)] there exist distinct vertices $u, v \in V(F) \setminus (A \cup B)$,
    \item[(iii)] there are at most $e(F)/4-k$ vertices of $F$ with degree more than 1. Moreover, every edge of $F$ contains at most one vertex of degree $1$, $F$ has no isolated vertices, and
    \item[(iv)] every $U\subseteq V(F)$ of size $|U|\geq 2$ satisfies $\Delta(U)\geq 2$.
\end{itemize}
\end{definition}

\noindent
Note that all eligible hypergraphs have $e(F)\geq 4k$, by (iii). The following lemma shows how to construct the hypergraph $F_{j+1}$ based on $F_j$ and shows that the eligibility is maintained under this construction. The hypergraphs $F_j$ constructed by repeated applications of this lemma will have deficiency $\Delta(F_j)=j$ and $e(F_j)\approx 2^j$ edges.

\begin{lemma}\label{lemma:maintaining eligibility}
Let $F$ be an eligible hypergraph with $\Delta(F)=k$. Suppose $F'$ is the hypergraph obtained by taking two vertex-disjoint copies of $F$ and only identifying the vertices of the sets $A$ in both copies. Then $F'$ is an eligible hypergraph with $\Delta(F')=k+1$ and $e(F')=2e(F)$.
\end{lemma}
\begin{proof}
Let us denote the two copies of $F$ by $F_1$ and $F_2$. Since $A$ is an independent set, no edges of $F_1$ and $F_2$ are identified and hence we have $e(F')=2e(F)$. Moreover, since exactly $k-1$ vertices are identified, we have $v(F')=2v(F)-(k-1)$ and therefore \[\Delta(F')=v(F')-e(F')=2v(F)-(k-1)-2e(F)=2k-(k-1)=k+1.\]

\begin{figure}[h!]
    \centering
    \begin{tikzpicture}
        % Define coordinates
  \def\R{1.5}

  % Draw semicircle
  \draw (-2*\R, \R) arc(90:270:\R);
  \draw (-2*\R, \R) -- (0, \R);
  \draw (-2*\R, -\R) -- (0, -\R);
  \draw (0, \R) arc(90:-90:\R);
  \draw[blue] (2*\R, \R) arc(90:-90:\R);
  \draw[blue] (0, \R) arc(90:270:\R);
  \draw[blue] (2*\R, \R) -- (0, \R);
  \draw[blue] (2*\R, -\R) -- (0, -\R);
  
  \draw (-1.5*\R, 0) circle(0.5);
  \coordinate[vtx] (v1) at (-2.5*\R, 0);
  \coordinate[vtx] (u1) at (-2.5*\R, -0.9);

  \draw[blue] (2.5*\R, 0) circle(0.5);
  \coordinate[vtx, blue] (v2) at (1.5*\R, 0);
  \coordinate[vtx, blue] (u2) at (1.5*\R, -0.9);

  \draw[red] (-1.5*\R, 0.6) arc(90:270:0.6);
  \draw[red] (-1.5*\R, 0.6) -- (1.5*\R, 0.6);
  \draw[red] (-1.5*\R, -0.6) -- (1.5*\R, -0.6);
  \draw[red] (1.5*\R, 0.6) arc(90:-90:0.6);
  % Annotations
  \node at (-1.2*\R, 1.2*\R) {$F_1$}; 
  \node[blue] at (1.2*\R, 1.2*\R) {$F_2$}; 
  \node[right] at (v1) {$v_1$};
  \node[right] at (u1) {$u_1$};
  \node at (-1.5*\R, 0) {$B_1$};
  \node[left, blue] at (v2) {$v_2$}; 
  \node[left, blue] at (u2) {$u_2$}; 
  \node[blue] at (2.5*\R, 0) {$B_2$};
  \node[red] at (0, 0.85) {$U'$}; 
  \node at (0, -0.8*\R) {$A$}; 
    \end{tikzpicture}
    \caption{Illustration of the proof of Lemma~\ref{lemma:maintaining eligibility}.}
\end{figure}

It remains to show that $F'$ is an eligible hypergraph. We begin by introducing some notation and defining the sets $A', B'$ and vertices $u', v'$ satisfying the conditions $\mathrm{(i)}$ and $\mathrm{(ii)}$ of Definition~\ref{defn:eligible}.

For $i\in \{1,2\}$, since $F_i$ is eligible, it contains a good set $B_i$ disjoint from $A$, and additional vertices $u_i, v_i \in V(F_i) \setminus (A\cup B_i)$. To show that $F'$ is eligible we define two sets $A' \coloneqq B_1\cup \{v_2\}$, $B'\coloneqq B_2\cup \{v_1\}$, along with additional vertices $u' \coloneqq u_1$ and $v' \coloneqq u_2$. Note that $u'$ and $v'$ are distinct and $u', v' \in V(F') \setminus (A' \cup B')$, verifying the second condition in Definition \ref{defn:eligible}. Finally, for every set $U'\subset V(F')$, we define the sets $U^{(1)}=U'\cap V(F_1)$ and $U^{(2)}=U'\cap V(F_2)$.

Let us now show that $A', B'\subset V(F')$ are good sets with respect to $F'$. Since the definitions of $A'$ and $B'$ are completely symmetric, it suffices to verify the conditions for only one of these sets, say $A'$. Note that $|A'|=|B_1|+1=k=\Delta(F')-1$, which shows that $A'$ has the required size. Further, $A'$ is an independent set since there are no edges between $F_1\backslash A$ and $F_2\backslash A$, and so, in particular, there are no edges between $B_1$ and $v_2$. To show that $A'$ is a good set we need to show that for any set $U'$ with $A'\subsetneq U'\subseteq V(F')$ one has $\Delta(U')\geq |A'|+1=k+1$. We will establish this inequality by using the properties of good sets $B_1\subseteq V(F_1)$ and $A\subseteq V(F_1)$. We have two cases: either $U^{(1)}\supsetneq B_1$ or $U^{(1)}=B_1$ (since $A'=B_1\cup \{v_2\}\subset U'$, we must have $B_1\subseteq U^{(1)}$). 

In the case that $U^{(1)}\supsetneq B_1$, we use the assumption that $B_1$ is a good set in $F_1$ to conclude $\Delta(U^{(1)})\geq |B_1|+1$. Furthermore, we know that $A\cap U^{(2)}$ is a good set in $F_2$ and $A\cap U^{(2)}\subsetneq U^{(2)}$ since $v_2\in U^{(2)}$. Thus, we have $\Delta(U^{(2)})\geq |A\cap U^{(2)}|+1 = |A\cap U'|+1$. These two bounds on $\Delta(U^{(1)})$ and $\Delta(U^{(2)})$ are sufficient to bound the deficiency of $U'$ in $F'$ as follows.
\begin{align*}
    \Delta(U')&=v(F'[U'])-e(F'[U']) \\ &=v(F_1[U^{(1)}])+v(F_2[U^{(2)}])-|U^{(1)}\cap U^{(2)}|-e(F_1[U^{(1)}])-e(F_2[U^{(2)}])\\
    &=\Delta(U^{(1)})+\Delta(U^{(2)})-|A\cap U'|\geq|B_1|+1+|A\cap U'|+1-|A\cap U'|=|B_1|+2=k+1.
\end{align*}

In the case that $U^{(1)}=B_1$, there are no vertices in $U^{(1)}\cap U^{(2)}=U^{(1)}\cap A$ and there are at least two vertices in $U^{(2)}$ since $B_1\cup \{v_2\}\subsetneq U'$. Hence, $\Delta(U^{(1)})=|B_1|$ since $B_1$ is an independent set, and $\Delta(U^{(2)})\geq 2$ since $F_2$ satisfies condition $\mathrm{(iv)}$ of Definition~\ref{defn:eligible}. Hence,
\begin{align*} \Delta(U') &=v(F_1[U^{(1)}])+v(F_2[U^{(2)}])-e(F_1[U^{(1)}])-e(F_2[U^{(2)}])\\
&=\Delta(U^{(1)})+\Delta(U^{(2)})\geq |B_1|+2=k+1.
\end{align*}

This completes the proof that $A'$ is a good set, thus verifying the first condition (of Definition~\ref{defn:eligible}) for $F'$ to be eligible. Further, to show that $F'$ satisfies the third condition, note that any vertex of $F'$ having degree more than 1 is either in $A$ or it already had degree more than 1 in $F_1$ or $F_2$. Hence, there are at most $2(e(F)/4-k)+k-1=e(F')/4-k-1$ vertices of degree more than 1 in $F'$. Moreover, it is clear that every edge of $F'$ has at most one vertex of degree $1$, since any vertex with degree $1$ in $F'$ must also have degree $1$ in $F_1$ or $F_2$. It is also clear that $F'$ has no isolated vertex, since $F_1$ and $F_2$ do not have an isolated vertex.

Finally, to show that $F'$ satisfies the fourth condition, we now verify that every set $U'\subseteq V(F')$ of size $|U'|\geq 2$ satisfies $\Delta(U')\geq 2$. Suppose for a contradiction that there is a set $U'\subseteq V(F')$ of size $|U'|\geq 2$ such that $\Delta(U')\leq 1$. Then, the assumption that $F_1$ and $F_2$ are eligible implies that we cannot have $U^{(1)}=U'$ or $U^{(2)}=U'$. This means that both $U^{(1)}$ and $U^{(2)}$ contain vertices outside $A$. Since $A\cap U'$ is a good set for both $F_1$ and $F_2$, we know that $\Delta(U^{(1)})\geq |A\cap U'|+1$ and $\Delta(U^{(2)})\geq |A\cap U'|+1$. Thus, 
\begin{align*}
    \Delta(U')&=v(F_1[U^{(1)}])+v(F_2[U^{(2)}])-|U^{(1)}\cap U^{(2)}|-e(F_1[U^{(1)}])-e(F_2[U^{(2)}])\\
    &=\Delta(U^{(1)})+\Delta(U^{(2)})-|A\cap U'|\geq |A\cap U'|+2\geq 2,
\end{align*}
completing the proof of the lemma.
\end{proof}

\noindent
To start the iteration, we need to show how to find eligible graphs. The following lemma shows how to construct an eligible hypergraph $F_0$ and find many copies of it in $\cH$.

\begin{lemma}\label{lemma:start}
Let $1\leq s\leq t$ be positive integers. Let $K_{s, t}^+$ be the hypergraph obtained by adding a single vertex $v_e$ to every edge $e\in E(K_{s, t})$ of the complete bipartite graph $K_{s, t}$. For any $0<\eps<1/s^3t$ and sufficiently large $n$, every linear hypergraph $\cH$ with $n$ vertices and $\Omega(n^{2-\eps})$ edges contains $\Omega(n^{s+t-s^3t\eps})$ copies of $K_{s, t}^+$. Moreover, if $8(s+t)\leq st$, the hypergraph $K_{s, t}^+$ is eligible and satisfies $\Delta(K_{s, t}^+)=s+t$ and $e(K_{s, t}^+)=st$.
\end{lemma}
\begin{proof}
To show that every linear hypergraph $\cH$ with $\Omega(n^{2-\eps})$ edges contains many copies of $K_{s, t}^+$, we take a tripartite subgraph $\cH'\subseteq \cH$ containing $\Omega(e(\cH))$ edges and whose parts all have size at least $n/4$ (which can be done by randomly partitioning the vertex set). Let us denote the three parts by $X, Y, Z$, and consider an auxiliary colored graph $G$ on the vertex set $X\cup Y$. Vertices $x\in X, y\in Y$ are adjacent in $G$ precisely when there exists a vertex $z\in Z$ such that $xyz$ is an edge of $\cH$. In this case, the edge $xy$ receives color $z$. Since $\cH$ is linear, each edge of $G$ is assigned a unique color. Furthermore, the coloring is proper for the same reason. Finally, $G$ contains $\Omega(n^{2-\eps})$ edges and $\Theta(n)$ vertices.

The reason behind defining the graph $G$ is the simple observation that every rainbow copy of $K_{s, t}$ in $G$ corresponds to a copy of $K_{s, t}^+$ in $\cH$. Hence, it suffices to show that there are $\Omega(n^{s+t-s^3t\eps})$ rainbow copies of $K_{s, t}$ in $G$. To do this, we use a simple lemma from the paper of Keevash, Mubayi, Sudakov and Verstra\"ete \cite{KMSV07} which states that in any properly colored copy of $K_{s, t'}$, for $t'>(s(s-1)+1)(t-1)$ one can find a rainbow copy of $K_{s, t}$. Let us fix $t'=s^2t$ and use the supersaturation of $K_{s, t'}$ in $G$ to obtain at least $\Omega(n^{s+t'-st'\eps})$ copies of $K_{s, t'}$ in $G$ (see e.g., the K\H{o}v\'{a}ri-S\'{o}s-Tur\'{a}n theorem~\cite{kHovari1954problem}). Then, using Lemma 2.3 from \cite{KMSV07}, in each of these copies one can find a rainbow copy of $K_{s, t}$. Since each copy of $K_{s, t}$ is contained in at most $n^{t'-t}$ copies of $K_{s, t'}$, we conclude that we have at least $\Omega(\frac{n^{s+t'-st'\eps}}{n^{t'-t}})=\Omega(n^{s+t-st'\eps})=\Omega(n^{s+t-s^3t\eps})$ rainbow copies of $K_{s, t}$, as desired.

The number of edges of $K_{s, t}^+$ is equal to the number of edges of $K_{s, t}$ and so $e(K_{s, t}^+)=st$. Further, since we add a new vertex to $K_{s, t}$ for every edge, the total number of vertices of $K_{s, t}^+$ is $v(K_{s, t}^+)=s+t+st$ and so $\Delta(K_{s, t}^+)=s+t$.

Finally, if $8(s+t)\leq st$, then we will show that $K_{s, t}^+$ is eligible. Note that if $T$ is a spanning tree of $K_{s, t}$, then the set $A_T=\{v_e | e\in T\}$ is a good set of size $s+t-1=\Delta(K_{s, t}^+)-1$. To verify this observation, let us take an arbitrary set $U\supsetneq A_T$ and show that $\Delta(U)\geq |A_T|+1$. First, observe that the statement is trivial if $U\cap V(K_{s,t})=\emptyset$ since in this case $U$ is an independent set and $\Delta(U)=|U|$. Hence, assume that $U\cap V(K_{s,t})\neq \emptyset$. Now note that it suffices to show $\Delta(U)\geq |A_T|+1$ only for sets $U$ completely contained within $A_T\cup V(K_{s, t})$. Once the statement is shown for such sets $U$, adding vertices outside $A_T\cup V(K_{s, t})$, which all have degree exactly 1 in $K_{s, t}^+$, cannot decrease the deficiency of $U$. Hence, we focus on sets $U$ with $A_T\subsetneq U\subseteq A_T\cup V(K_{s, t})$. The number of edges of $K_{s, t}^+$ spanned by such a set $U$ is precisely the number of edges of $T$ spanned by the intersection $U\cap V(K_{s, t})$, i.e. $e(K_{s, t}^+[U])=e(T[U\cap V(K_{s, t})])$. Since $T$ is a tree, we must have $e(T[U\cap V(K_{s, t})])\leq |U\cap V(K_{s, t})|-1$ and so 
\[\Delta(U)=|A_T|+|U\cap V(K_{s, t})|-e(K_{s, t}^+[U])\geq |A_T|+|U\cap V(K_{s, t})|-|U\cap V(K_{s, t})|+1=|A_T|+1.\]

Hence, to construct two disjoint good sets in $K_{s, t}^+$ of size $\Delta(K_{s, t}^+)-1=s+t-1$, it suffices to take two edge-disjoint spanning trees of $K_{s, t}$, which is certainly possible as the condition $st\geq 8(s+t)$ implies $s, t\geq 9$ (verifying (i) of Definition~\ref{defn:eligible}). Of course, $K_{s, t}^+$ contains two distinct additional vertices besides these good sets (verifying (ii) of Definition~\ref{defn:eligible}). Further, the vertices of degree more than $1$ in $K_{s, t}^+$ are precisely the initial vertices of $K_{s, t}$ and we have at most $s+t\leq st/4-\Delta(K_{s, t}^+)$ such vertices by the assumption that $st\geq 8(s+t)$. Moreover, every edge of $K_{s, t}^+$ contains exactly one vertex of degree 1 and $K_{s,t}^+$ has no isolated vertex (verifying (iii) of Definition~\ref{defn:eligible}). Finally, since every edge of $K_{s, t}^+$ contains a distinct vertex of degree $1$ and two vertices of $V(K_{s, t})$, any set of $e$ edges of $K_{s, t}^+$ must contain at least $e+2$ distinct vertices. This shows that there are no sets $U\subset V(K_{s, t}^+)$ of deficiency $\Delta(U)\leq 1$ and size $|U|\geq 2$ (verifying (iv) of Definition~\ref{defn:eligible}). Hence, $K_{s, t}^+$ is eligible.
\end{proof}

\subsection{Iteration}\label{subsec:iteration}

We begin this section by showing that if we have many copies of a smaller hypergraph $F_j$ in $\cH$, then we can either glue them together to produce many copies of a larger hypergraph $F_{j+1}$ or  we find an alternative structure in $\cH$ (as mentioned in the proof overview). We now define this structure.

\begin{definition}
Let $F$ be a hypergraph and let $r$ be a positive integer. A hypergraph $\tilde F$ is an \textit{$(r, F)$-sunflower} if there exists a set $U\subsetneq V(F)$ with $\Delta(U)\geq \Delta(F)$ and $r$ embeddings $\varphi_1, \dots, \varphi_r:V(F)\to V(\tilde F)$ satisfying the following two conditions:
\begin{itemize}
    \item[(i)] the embeddings $\varphi_1, \dots, \varphi_r$ cover $\tilde F$, i.e. $\bigcup_{i=1}^r \varphi_i(V(F))=V(\tilde F)$, and 
    \item[(ii)] for different indices $1\leq i<j\leq r$ we have $\varphi_i(u)=\varphi_j(v)$ if and only if $u=v\in U$. In particular, the sets $\varphi_i(V(F)\backslash U)$ and $\varphi_j(V(F)\backslash U)$ are disjoint when $i\neq j$.
\end{itemize}
\end{definition}

The main reason sunflowers are useful for us is that the deficiency of $\tilde F$ is bounded by the deficiency of $F$. Indeed, $v(\tilde F)=|U|+r(v(F)-|U|)$ and $e(\tilde F)=e(F[U])+r(e(F)-e(F[U]))$, so the deficiency of $\tilde F$ can be bounded as follows. $\Delta(\tilde F)=v(\tilde F)-e(\tilde F)=\Delta(F)+(r-1)(\Delta(F)-\Delta(U))\leq \Delta(F)$.

The following lemma, whose proof uses ideas from the proof of Lemma 3.1 in \cite{CGLS23}, shows that any hypergraph $\cH$ with many copies of $F$ contains either an $(r, F)$-sunflower or many copies of $F'$ (obtained by gluing copies of $F$). In the proof of Theorem~\ref{thm:power saving} we will mostly use this lemma with $m=2$.

\begin{lemma}\label{lemma:iteration step}
Let $r, m\geq 2$ be fixed integers and let $F$ be an eligible hypergraph with deficiency $\Delta(F)=k$. Then, there exists a hypergraph $F'$ satisfying $\Delta(F')=k+m-1$, $e(F')=m\cdot  e(F)$ and a positive constant $n_0=n_0(r,m,F)$ with the following property: for all $\eps\leq 1/2$ and all hypergraphs $\cH$ on $n>n_0$ vertices containing $\Omega(n^{k-\eps})$ copies of $F$, either
\begin{itemize}
    \item $\cH$ contains an $(r, F)$-sunflower, or
    \item $\cH$ contains $\Omega(n^{k+m-1-m\eps})$ copies of $F'$.
\end{itemize}
Moreover, if $m=2$, the resulting graph $F'$ is also eligible.
\end{lemma}
\begin{proof}
We take $F'$ to be the hypergraph obtained by gluing $m$ copies of $F$ along a good set $A\subset V(F)$ of size $k-1$. In other words, to obtain $F$ we start from $m$ disjoint copies $F^{(1)}, \dots, F^{(m)}$ and for each vertex $v\in A$ we identify the $m$ instances of that vertex with each other. For $m=2$, this construction is identical to the construction from Lemma~\ref{lemma:maintaining eligibility} and therefore the resulting hypergraph $F'$ is indeed eligible. 

We now focus on proving the main statement of the lemma. Throughout the proof, we identify the vertex set of $F$ with $[v(F)]$, where the set $A \subset V(F)$ is identified with $[k-1]$, and we think of copies of $F$ in $\cH$ as embeddings $\varphi:V(F)\to V(\cH)$. As a first step, we randomly partition vertices of $\cH$ into $v(F)$ parts $S_1, \dots, S_{v(F)}$. We call a copy of $F$ \textit{proper} if the corresponding embedding $\varphi:[v(F)]\to V(\cH)$ satisfies $\varphi(i)\in S_i$ for all $i\in V(F)$. The expected number of proper copies of $F$ is $v(F)^{-v(F)}$ times the total number of copies of $F$ in $\cH$. Hence, we can fix a partition $S_1, \dots, S_{v(F)}$ with $\Omega(n^{k-\eps})$ proper copies of $F$. Let us denote the collection of these copies by $\cF$. For two embeddings $\varphi_1, \varphi_2\in \cF$, we define their intersection $U(\varphi_1, \varphi_2)=\{i\in V(F):\varphi_1(i)=\varphi_2(i)\}$. The random partitioning we performed ensures that any overlap between two copies $\varphi_1, \varphi_2\in \cF$ can only happen on the vertices playing the same role in both copies, i.e. we can never have $\varphi_1(i)=\varphi_2(j)$ for $i\neq j$. Formally, for any $\varphi_1, \varphi_2\in \cF$ we have $\varphi_1(V(F))\cap \varphi_2(V(F))=\varphi_1(U(\varphi_1, \varphi_2))$.

We now define an auxiliary graph $\cG$ whose set of vertices is $\cF$. We define two copies $\varphi_1, \varphi_2\in \cF$ to be adjacent in $\cG$ if $A\subsetneq U(\varphi_1, \varphi_2)$. We now have two cases - either the auxiliary graph $\cG$ has a vertex of degree at least $v(F)!(2r)^{v(F)}$ or all vertices of $\cG$ have degree less than $v(F)!(2r)^{v(F)}$.

Suppose we are in the first case and we have an embedding $\varphi_0:V(F)\to V(\cH)$ with at least $v(F)!(2r)^{v(F)}$ neighbours in $\cG$. Then, by the pigeonhole principle one can find a set $U_0$ such that at least $v(F)!r^{v(F)}$ embeddings $\varphi\in N_{\cG}(\varphi_0)$ satisfy $U(\varphi, \varphi_0)=U_0$. We denote the set of these embeddings by $\cF_0$. Let us now recall the Erd\H{o}s-Rado sunflower lemma, which we plan to apply to the vertex sets of the copies of $F$ in $\cF_0$.

\begin{theorem*} [Erd\H{o}s-Rado sunflower lemma, \cite{ER60}]

Let $\cS$ be a family of sets, where each set has cardinality $k$. If $|\cS|\geq k!(r-1)^k$, there exist $r$ distinct sets $S_1, \dots, S_r\in \cS$ which satisfy $S_i\cap S_j=S_1\cap \dots\cap S_r$ for any $1\leq i< j\leq r$.
\end{theorem*}

\noindent
We apply the Erd\H{o}s-Rado sunflower lemma to the family $\cS=\{\varphi(V(F))|\varphi\in \cF_0\}$. Note that each set in this family has cardinality $v(F)$. Since $|\cS|\geq v(F)!r^{v(F)}$, we obtain a set of $r$ different embeddings $\varphi_1, \dots, \varphi_r$ and a set $U\subsetneq V(F)$ such that $U(\varphi_i, \varphi_{j})=U$ for all $i, j\in [r]$, $i\neq j$. We claim that the union of these $r$ copies of $F$ forms an $(r, F)$-sunflower contained in $\cH$. Since $U(\varphi_0, \varphi_i)\cap U(\varphi_0, \varphi_{j})\subseteq U(\varphi_i, \varphi_{j})$ for all $i, j\in [r]$, we see that $U_0\subseteq U$. By the definition of the adjacency relation in $\cG$, we have $A\subsetneq U_0\subset U$ and therefore the set $U$ has deficiency $\Delta(U)\geq |A|+1=\Delta(F)$. Hence,  $\bigcup_{i=1}^r \varphi_i(F)$ really is an $(r, F)$-sunflower.

Now, we consider the second case, in which all vertices of $\cG$ have degree less than $v(F)!(2r)^{v(F)}$. In this case, $\cG$ contains an independent set $I$ of size at least $\frac{|V(\cG)|}{v(F)!(2r)^{v(F)}}=\Omega(n^{k-\eps})$. For each $(k-1)$-tuple $\mathbf{v}=(v_1, \dots, v_{k-1})\in S_1\times\dots\times S_{k-1}$ representing the vertices of $A$, we define $\cF_{\mathbf{v}}$ to be the set of embeddings $\varphi\in I$ having $\varphi(i)=v_i$ for all $i\in [k-1]$. Note that any two embeddings $\varphi_1, \varphi_2\in \cF_{\mathbf{v}}$ have $U(\varphi_1, \varphi_2)=A=[k-1]$ since they are not adjacent in the graph $\cG$. Therefore, for any $m$ distinct embeddings $\varphi_1, \varphi_2, \ldots, \varphi_m \in \cF_{\mathbf{v}}$, the subgraph $\varphi_1(F)\cup \ldots \cup \varphi_m(F)$ of $\cH$ is a copy of $F'$. In particular, this shows that there are $\Omega\left(\binom{|\cF_{\mathbf{v}}|}{m}\right)$ copies of $F'$ containing the $(k-1)$-tuple $\mathbf{v}$. Since the copies of $F'$ corresponding to different $(k-1)$-tuples $\mathbf{v}$ are different, the number of copies of $F'$ in $\cH$ is at least $\sum_{\substack{\mathbf{v}\in S_1\times\cdots\times S_{k-1}\\|\cF_{\mathbf{v}}|\geq m}} \binom{|\cF_{\mathbf{v}}|}{m}$. Since $x\mapsto \binom{x}{m}$ is a convex function when $x\geq m$, applying Jensen's inequality gives
\begin{align*}
    \sum_{\substack{\mathbf{v}\in S_1\times\cdots\times S_{k-1}\\|\cF_{\mathbf{v}}|\geq m}} \binom{|\cF_{\mathbf{v}}|}{m}&\geq \prod_{i=1}^{k-1}|S_i|\binom{(\prod_{i=1}^{k-1}|S_i|)^{-1}\sum_{|\cF_{\mathbf{v}}|\geq m} |\cF_{\mathbf{v}}|}{m} \\&\geq \frac{\left(\sum_{\mathbf{v}} |\cF_{\mathbf{v}}|-mn^{k-1}\right)^m}{m^m (\prod_{i=1}^{k-1}|S_i|)^{m-1}}
    \geq \Omega\left(\frac{n^{mk-m\eps}}{n^{(m-1)(k-1)}}\right)= \Omega\left(n^{m+k-1-m\eps}\right),
\end{align*}
where the last inequality comes from the fact that the sum $\sum_{\mathbf{v}} |\cF_{\mathbf{v}}|$ simply counts the number of proper copies of $F$ in $\cH$ and that each of the sets $S_1, S_2, \ldots, S_k$ has size at most $n$. Therefore, $\cH$ contains $\Omega(n^{k+m-1-m\eps})$ copies of $F'$, as needed. This completes the proof of the lemma.
\end{proof}


\begin{lemma}\label{lemma:sunflower}
Suppose that $F$ is an eligible hypergraph and suppose that a hypergraph $\cH$ contains an $(r, F)$-sunflower (for some $r$) with at least $e$ edges. If $e(F)\leq e/2$, then $\cH$ contains an $(e+\Delta(F), e)$-configuration.
\end{lemma}
\begin{proof}
By our assumption, $\cH$ contains an $(r, F)$-sunflower with at least $e$ edges. Since it is a sunflower, it has deficiency at most $\Delta(F)$. However, this sunflower may have too many edges and vertices, so it does not immediately give us an $(e+\Delta(F), e)$ configuration. Thus, the idea is to remove some vertices of degree $1$ from the copies of $F$ forming this sunflower and thus reduce the number of edges and vertices of the sunflower, without changing the deficiency in the process.

Suppose the copies of $F$ forming this sunflower correspond to embeddings $\varphi_1, \dots, \varphi_r$ and suppose that these copies are glued over a set $U\subset V(F)$ with $\Delta(U)\geq \Delta(F)$. We call the set $U$ the \textit{core} and  for each copy of $F$ forming this sunflower, we call the set $P=V(F)\backslash U$ a \textit{petal}. Further, let $e_U=e(F[U])$ denote the number of edges of $F$ completely contained in the core and let $e_P=e(F)-e(F[U])$ denote the number of edges incident to a vertex in the petal $P$. The assumption that $\Delta(U)\geq \Delta(F)$ implies that $e_P\geq |P|$, which means that every petal adds at least as many edges as vertices. Note that the sunflower constructed using only the copies $\varphi_1, \dots, \varphi_p$ would have $p|P|+|U|$ vertices and $pe_P+e_U$ edges. Let us fix $p$ to be the smallest number of petals one needs to attach to the core to get at least $e$ edges in the sunflower, i.e. $p=\lceil\frac{e-e_U}{e_P}\rceil$. Note that we may assume that $p\geq 3$ since $p=2$ implies that the sunflower has exactly $e$ edges (because we assumed $e(F) \le e/2$) and hence it represents an $( e+\Delta(F), e)$-configuration.

The total number of vertices in this sunflower is $p|P|+|U|$ and thus if $p|P|+|U|\leq e+\Delta(F)$, we are done since we have found an $(e+\Delta(F), e)$-configuration. Otherwise, $p|P|+|U|> e+\Delta(F)$, in which case our goal is to remove some vertices of degree 1 to obtain exactly $e+\Delta(F)$ vertices. Then, since the deficiency is still at most $\Delta(F)$, this gives us the desired $(e+\Delta(F), e)$-configuration. We claim that it is always sufficient to remove at most $|P|$ vertices of degree $1$ in order to bring the number of vertices to $e+\Delta(F)$. Indeed, note that although the number of vertices of the sunflower formed by $\varphi_1, \dots, \varphi_p$ is larger than $e+\Delta(F)$, it cannot be much larger. Namely, by the minimality of $p$ we have $(p-1) e_P+e_U\leq e$ and so 
\begin{align*}
    p|P|+|U|&= (p-1)|P|+ v(F)\leq (p-2)e_P + |P| + e(F)+\Delta(F)\\
    &= (p-1)e_P+e_U+|P|+\Delta(F)\leq e+\Delta(F)+|P|.
\end{align*}

Thus, it suffices to show that the sunflower consisting of $\varphi_1, \dots, \varphi_p$ has at least $|P|$ vertices of degree $1$. Let $v_U$ denote the number of vertices $v$ of $U$ such that $v$ has degree $1$ in $F$ and the edge incident to $v$ does not contain any vertices of $P$, and let $v_P$ be the number of vertices of $P$ which have degree $1$. Then the number of vertices of degree $1$ in the sunflower consisting of $\varphi_1,\dots,\varphi_p$ is equal to $v_U+pv_P$. Our goal is to show that $v_U+pv_P\geq |P|$. Suppose this is not the case and we have $v_U+pv_P<|P|$. Since $F$ is eligible, at most $e(F)/4$ vertices in the hypergraph $F$ can have degree more than $1$, so we have $v_P\geq |P|-e(F)/4$. We know that $p\geq 3$, so $v_U <  |P|-3v_P$.

Since $F$ is eligible, at least $3e(F)/4$ vertices in $F$ have degree $1$ and moreover, every edge of $F$ contains at most one vertex of degree $1$. So there are at least $3e(F)/4$ edges in $F$ incident to a vertex of degree $1$. Out of these, only $v_U$ are completely contained in $U$. Thus, the number of edges incident to $P$ is at least \[e_P\geq 3e(F)/4-v_U\geq 3e(F)/4-|P|+3v_P\geq 3e(F)/4-|P|+3(|P|-e(F)/4)=2|P|.\] 

In other words, every petal adds twice as many edges as vertices. We will show that this contradicts our assumption that $p|P| + |U| > e + \Delta(F)$. More precisely, we have 
\begin{align*}
p|P|+|U|\leq &\left(\frac{e-e_U}{e_P}+1\right)|P| +|U|\leq \frac{e}{e_P}|P|+|P|+|U|\\
\leq & \frac{e}{e_P}|P|+v(F)\leq \frac{e}{2}+e(F)+\Delta(F)\leq e+\Delta(F),
\end{align*}
since $e(F) \le e/2$. This completes the proof of the lemma.
\end{proof}

Now we have all the necessary ingredients to prove Theorem~\ref{thm:power saving}.

\begin{proof}[Proof of Theorem~\ref{thm:power saving}.]
Throughout the proof, we will assume that we have a linear hypergraph $\cH$ with $\Omega(n^{2-\eps})$ edges, where $\eps<e^{-5}$ and $e$ is the required number of edges in the final configuration. Note that the assumption that $\cH$ is linear is without any loss of generality, since if any of the pairs $(x, y)\in V(\cH)^2$ was in at least $e$ edges then we would have an $(e+2, e)$-configuration. Hence, one can keep a positive constant (dependent only on $e$) fraction of all edges in $\cH$ while ensuring that no pair $(x, y)\in V(\cH)^2$ appears in more than one edge.

In case $e$ is small, say $e<512$, we can use Lemma~\ref{lemma:start} to find many copies of $K_{\lceil \sqrt{e}\rceil, \lceil \sqrt{e}\rceil}^{+}$ in $\cH$. The deficiency of this hypergraph is $2\lceil\sqrt{e}\rceil\leq \lfloor \log_2 e\rfloor +38$ as long as $e<512$, which is sufficient for the proof of our theorem in this case.

For $e\geq 512$, we apply Lemma~\ref{lemma:start} with $s=t=16$ to find an eligible hypergraph $F_0$ such that $\cH$ contains $\Omega(n^{k_0-s^3 t \eps})$ copies of $F_0$, where $k_0=\Delta(F_0) = 32$. Note that $e(F_0) = 256$. Let us now fix the parameters $r=e$ and $\ell=\lfloor \log_2 e/e(F_0)\rfloor-1$ (where $\ell\geq 0$ since $e\geq 512$). Observe that the parameters were chosen so that $2^\ell s^3t \eps\leq 1/2$. Having fixed these parameters, we will show how to find a $(v, e)$-configuration in $\cH$ with $v-e\leq \lfloor \log_2 e\rfloor +26$.

The idea is to iteratively apply Lemma~\ref{lemma:iteration step} with $m=2$ to construct a sequence of eligible hypergraphs $F_0, F_1, F_2, \dots, F_\ell$ with the property that either $\cH$ contains an $(r, F_j)$-sunflower for some $0\leq j\leq \ell$ or $\cH$ contains $\Omega(n^{k_0+j-2^j \eps'})$ copies of $F_j$ for each hypergraph $F_j$ in this sequence, where $\eps'=s^3t\eps$. Indeed, the construction of this sequence is given by Lemma~\ref{lemma:iteration step}, where we set $F_{j+1}=F_j'$ for every $0 \le j \le \ell-1$. Since $e(F_j')=2e(F_j)$ and $\Delta(F_j')=\Delta(F_j)+1$, this construction satisfies $e(F_j)=2^je(F_0)$ and $\Delta(F_j)=\Delta(F_0)+j$ for all $0 \le j \le \ell$.
If we assume that $\cH$ does not contain an $(r, F_j)$-sunflower for any $0\leq j\leq \ell$, one can show by induction on $j$ that $\cH$ contains $\Omega(n^{k_0+j-2^j \eps'})$ copies of $F_j$ for each $0\leq j\leq \ell$. The base case is given by Lemma~\ref{lemma:start}, since $\cH$ contains $\Omega(n^{k_0-\eps'})$ copies of $F_0$. For $j\ge 1$, from the assumption that $\cH$ contains $\Omega(n^{k_0+j-1-2^{j-1}\eps'})$ copies of $F_{j-1}$ and that $\cH$ contains no $(r, F_{j-1})$-sunflower, Lemma~\ref{lemma:iteration step} allows us to deduce that $\cH$ contains $\Omega(n^{k_0+j-2^j \eps'})$ copies of $F_j$, allowing us to continue the induction. Recall that the parameters were set such that $2^j \eps'\leq 2^\ell s^3 t\eps \leq \frac{1}{2}$ and hence Lemma~\ref{lemma:iteration step} can be applied for all $j\leq \ell$.

If $\cH$ contains an $(r, F_j)$-sunflower for some $0\leq j\leq \ell$, then Lemma~\ref{lemma:sunflower} shows that $\cH$ contains an $(e+\Delta(F_j), e)$-configuration. Indeed, the assumptions of Lemma~\ref{lemma:sunflower} are satisfied since $e(F_j)\leq 2^\ell e(F_0)\leq \frac{e}{2e(F_0)}e(F_0)=e/2$ and $r=e$ (which implies that the sunflower has at least $e$ edges, since $F_j$ has no isolated vertex). Moreover, note that finding an $(e+\Delta(F_j), e)$-configuration in $\cH$ would indeed suffice to prove Theorem~\ref{thm:power saving} since
\[\Delta(F_j)=\Delta(F_0)+j\leq 32+\ell\leq 32+\lfloor \log_2 e/e(F_0) \rfloor -1= \lfloor \log_2 e\rfloor +23.\]

The only case left to consider is when $\cH$ does not contain an $(r, F_j)$-sunflower for any $0\leq j\leq \ell$. Applying Lemma~\ref{lemma:iteration step} to $F_\ell$ with $m=4$ and using the fact we have no $(r, F_\ell)$-sunflower, we arrive at the conclusion that $\cH$ contains a hypergraph $F_\ell'$ obtained by gluing $4$ copies of $F_\ell$ along a good set $A\subset V(F_\ell)$ of size $\Delta(F_\ell)-1$. The hypergraph $F_\ell'$ satisfies $\Delta(F_\ell')=\Delta(F_\ell)+3$ and $e(F_\ell')=4e(F_\ell)\in [e, 2e]$. Since $F_\ell'$ is formed by gluing $4$ copies of $F_\ell$ along a set of size $\Delta(F_\ell)-1$ and $F_\ell$ is eligible, the number of vertices of degree more than $1$ in $F_\ell'$ is at most $4(e(F_\ell)/4-\Delta(F_\ell))+\Delta(F_\ell)-1\leq e(F_\ell')/4$. Hence, $F_\ell'$ contains at least $v(F_\ell')-e(F_\ell')/4 = e(F_\ell')+\Delta(F_\ell')-e(F_\ell')/4 \ge e(F_\ell')/2$ vertices of degree $1$, so one can remove $e(F_\ell')-e\leq e(F_\ell')/2$ vertices of degree $1$ from $F_\ell'$ without changing the deficiency of the configuration. In this way, for any $e \ge 512$, one also obtains an $(e+\Delta(F_\ell'), e)$-configuration with \[ \Delta(F_\ell')\leq \ell+3+\Delta(F_0)\leq 32+\lfloor \log_2 e/e(F_0)\rfloor+2= \lfloor \log_2 e\rfloor +26.\]
This completes the proof of Theorem~\ref{thm:power saving}.
\end{proof}
